\chapter{Summary and Conclusion}
\label{ch:conclusion}

\section{Summary of Contributions}
\label{sec:conclusion:summary}

This thesis developed methods for representing, generating, and evaluating
graphs. The work was motivated by a common difficulty across these tasks:
graphs are irregular, permutation-sensitive, and often constrained by
domain-specific validity requirements. A useful method must therefore preserve
graph structure, remain inspectable where possible, and support reliable
comparison when used inside a generative pipeline.
The thesis made three main contributions.

\textbf{NSPPK for interpretable graph representation.}
The first contribution was the Neighborhood Subgraph Pairwise Path Kernel
(NSPPK), an explicit graph-to-vector representation for node-attributed graphs. For
each graph, NSPPK constructs a sparse feature vector by enumerating structural
patterns around pairs of nodes. Each feature records the neighbourhoods around the nodes together with the shortest-path subgraph that connects them. In this way, the representation captures both local node-centred structure and the broader context linking different parts of the graph.
Continuous node attributes are incorporated directly into these explicit
features, avoiding discretisation while retaining deterministic, sparse, and
interpretable feature construction.

Empirically, NSPPK showed that carefully designed explicit representations can
remain competitive with learned graph neural models, especially in low-data or
sample-efficient regimes. Its sparse feature map also supports computationally
transparent downstream learning and enables predictive behaviour to be related
back to structural graph patterns.

\textbf{GCDG for graph-conditioned generation.}
The second contribution was the Graph-Conditioned Diffusion Generator (GCDG), a
conditional graph generator that uses graph-level vectors to guide the
generation of node-level representations. The central contribution is the
graph-level-to-node-level conditional generation formulation: a fixed
graph-level descriptor conditions the generation of a padded node-level feature
matrix, from which the final graph is decoded. The diffusion process is defined
over continuous node-level features, while the denoising network predicts the
clean node representation and latent states used by downstream prediction
modules. These modules estimate node presence, node labels, degree-related
information, edge probabilities, and edge labels where available. A
constraint-aware decoder then combines the node-level and pairwise predictions
to construct the final discrete graph.

This  design makes the generation process more inspectable than a single
black-box graph decoder. Intermediate quantities such as node-existence
probabilities, edge scores, degree predictions, and label predictions can be
examined before the final graph is selected. Experiments on synthetic
cycle--leg graphs and QM9 molecular graphs showed that conditioning quality and
graph validity are distinct issues: a model may produce plausible graphs while
failing to preserve the requested conditional structure. The chapter therefore
emphasised generation as conditional graph construction rather than
unconditional sampling alone.

\textbf{RankGen for statistically grounded evaluation.}
The third contribution was RankGen, a classifier-based framework for evaluating
generative models. RankGen treats evaluation as a collection of predictive
probes rather than as a single scalar similarity computation. The framework
defines four complementary metrics: Quality, Utility, Indistinguishability,
and Similarity, together with a composite score, Exchangeability. These metrics
diagnose different failure modes, including poor
task fidelity, lack of augmentation value, distributional separability, local
structural mismatch, and memorisation.

RankGen combines these probes with PAC-style finite-sample reliability
bounds and robust quantile-based aggregation. This separates the question of
whether a comparison is statistically reliable from the question of which
reliable model should be ranked higher. Experiments across controlled
perturbations, molecular graph
generators, image benchmarks, and open-domain text generation demonstrated that
the framework can expose behaviours that conventional metrics often collapse
or obscure.

\section{Key Insights and Lessons Learned}
\label{sec:conclusion:insights}

Several broader lessons emerge from these contributions.

First, explicit graph representations remain valuable. Although neural graph
models are highly expressive, deterministic feature maps can provide strong
performance, favourable sample efficiency, and clearer interpretability. NSPPK
shows that the limitations of classical kernels are not inherent to explicit
representations themselves; they can be addressed by enriching the structural
features and by treating continuous attributes directly.

Second, graph generation benefits from separating modelling roles. In GCDG, the
diffusion model is not asked to solve every part of graph generation at once.
Continuous structural representations, categorical node labels, edge
prediction, and final constraint satisfaction are handled by different
components. This separation does not remove the difficulty of graph generation,
but it makes the source of successes and failures easier to inspect.

Third, conditional generation must be evaluated conditionally. It is not enough
for a generated graph to look plausible under aggregate statistics. If the
generator receives a graph-level condition, then the generated graph should be
tested against the information encoded in that condition. The synthetic
cycle--leg experiments illustrate this point clearly: different models can fail
in different conditional regimes even when their outputs remain graph-like.

Fourth, generative model evaluation is inherently multi-dimensional. A
generator may preserve predictive signal but be distributionally
distinguishable; it may generate realistic-looking samples that add little
utility; or it may perform well under global metrics while failing local
mixing or memorisation checks. RankGen addresses this by treating evaluation as
a structured statistical problem rather than a search for one universal score.

\section{Limitations}
\label{sec:conclusion:limitations}

The work in this thesis also has limitations.

For NSPPK, the use of explicit substructure enumeration provides
interpretability but can become expensive on dense or very large graphs. The
method is designed around bounded radii and distances, so its practical
scalability depends on parameter choices and graph density. The feature hashing
pipeline is efficient and deterministic, but hashed representations can still
merge distinct patterns when the hash space is limited. In addition, NSPPK
captures structural information through predefined features rather than
learning task-specific representations end-to-end.

For GCDG, the  architecture introduces dependencies between components.
If the node representation does not preserve information needed for decoding,
the constrained decoder cannot recover it later. Conversely, a strong decoder or domain-specific constraint can improve validity,
but it also makes it harder to determine how much of the final performance
comes from the learned generator and how much comes from post-processing. 
than as a complete molecular-design benchmark.

For RankGen, the framework depends on the quality of the chosen classifiers,
feature representations, and resampling protocol. Classifier-based evaluation
is powerful because it connects generation to predictive tasks, but it also
inherits the biases and limitations of the classifiers and embeddings used.
The PAC-style bounds provide conservative reliability checks, yet they do not
eliminate all sources of experimental variability. RankGen should therefore be
viewed as a diagnostic framework rather than as a replacement for all
domain-specific evaluation.

Across the thesis, a further limitation is that the strongest claims are made
within the graph and benchmark settings studied here. Across the thesis, a further limitation concerns the scope of the evaluation
formalism. RankGen is developed primarily around supervised, classifier-based
probes, where predictive accuracy and real-versus-generated discrimination
provide natural diagnostic signals. Extending the framework to regression
tasks, unsupervised learning, or settings without reliable labels would require
additional formal development, including appropriate probe definitions,
uncertainty bounds, and criteria for task-relevant utility.

\section{Future Research Directions}
\label{sec:conclusion:future}

Several directions follow naturally from this work.

\textbf{Hybrid explicit and neural representations.}
NSPPK suggests that explicit graph features can provide useful structural
signals. A promising direction is to combine such features with neural graph
models, either as input channels, regularisers, attention biases, or
interpretable probes for learned embeddings. This could preserve the
sample-efficiency and auditability of kernels while benefiting from the
task-specific flexibility of neural representation learning.

\textbf{Scalable conditional graph generation.}
Future work on GCDG could improve scalability to larger and denser graphs by
developing more efficient decoding procedures, sparse pairwise prediction, and
hierarchical generation. Another direction is to learn constraint handling more
closely with the denoising process, reducing the gap between neural generation
and final feasibility projection while retaining inspectability.

\textbf{Richer condition preservation metrics.}
Conditional generation requires evaluation methods that compare generated
objects not only to a population, but also to the specific condition that
produced them. Future work could extend RankGen-style probes to conditional
settings by measuring condition-response fidelity, conditional diversity, and
the stability of generated samples under controlled changes to the condition
vector.

\textbf{Evaluation under privacy and safety constraints.}
RankGen already addresses memorisation as a diagnostic concern, but future work
could connect the framework more directly to privacy auditing, extraction
risk, and safety-sensitive deployment. This is especially important in domains
where generated data may contain sensitive information or where synthetic data
is used to support downstream decisions.

\textbf{Application-specific scientific validation.}
For molecular and other scientific graph domains, further validation should
include stronger domain-specific constraints, expert-defined utility criteria,
and prospective testing. Distributional similarity is useful, but scientific
generators ultimately need to support discovery, optimisation, or reliable
simulation in the target domain.

\section{Closing Remarks}
\label{sec:conclusion:closing}

The central message of this thesis is that effective and trustworthy graph
generation requires more than a powerful sampler. It requires representations
that expose relevant graph structure, generators whose intermediate decisions
can be inspected, and evaluation protocols that distinguish genuine utility
from surface-level similarity.

NSPPK, GCDG, and RankGen address these three parts of the problem. NSPPK shows
that explicit graph structure can be represented in a sparse and interpretable
form, including continuous attributes. GCDG shows how such graph-level
information can guide a staged denoising generator for discrete graphs. RankGen
shows how generated data can be assessed through statistically grounded,
diagnostic probes.

Taken together, the contributions support a broader view of graph generative
modelling: representation, generation, and evaluation should be designed
together. When these components are aligned, generated graphs can be assessed
not only by whether they appear plausible, but by whether they preserve
meaningful structure, satisfy relevant constraints, and provide reliable value
for downstream use.
